\documentclass[11pt]{article}
\usepackage{amsmath, amssymb, amsthm}
\usepackage{geometry}
\geometry{margin=1in}
\usepackage{algorithm}
\usepackage{algpseudocode}
\usepackage{color}

\begin{document}

\title{Convergence Analysis of Proximal Gradient Descent}
\author{}
\date{}
\maketitle

\section*{Algorithm Name}
Proximal Gradient Descent (PGD), also known as Forward-Backward Splitting.

\section*{Mathematical Formulation}
The algorithm aims to minimize a composite objective function of the form:
$$F(x) = f(x) + g(x)$$
where $f(x)$ is a differentiable function (typically representing a loss) and $g(x)$ is a potentially non-differentiable convex function (typically representing a regularizer). 

The update rule consists of a forward (gradient) step on $f$ followed by a backward (proximal) step on $g$:
$$y_k = x_k - \eta \nabla f(x_k)$$
$$x_{k+1} = \text{prox}_{\eta g}(y_k)$$
Combining these, the singular update rule is:
$$x_{k+1} = \text{prox}_{\eta g}(x_k - \eta \nabla f(x_k))$$
where the proximal operator is defined as:
$$\text{prox}_{\eta g}(y) = \arg\min_{x} \left\{ g(x) + \frac{1}{2\eta} \|x - y\|^2 \right\}$$

\section*{Pseudocode}
\begin{algorithm}[H]
\caption{Proximal Gradient Descent (PGD)}
\begin{algorithmic}[1]
\State \textbf{Input:} Initial point $x_0 \in \mathbb{R}^d$, step size $\eta > 0$, number of iterations $K$
\State \textbf{Initialize:} $k = 0$
\While{$k < K$}
    \State Compute gradient of smooth part: $d_k = \nabla f(x_k)$
    \State Take gradient descent step: $y_k = x_k - \eta d_k$
    \State Apply proximal operator: $x_{k+1} = \text{prox}_{\eta g}(y_k)$
    \State $k = k + 1$
\EndWhile
\State \textbf{Output:} $x_K$
\end{algorithmic}
\end{algorithm}

\section*{Dry Run Example}
Let $f(w) = (w-3)^2$ and $g(w) = |w|$ (L1 regularization). We set $w_0 = 0$ and $\eta = 0.1$.
The gradient is $\nabla f(w) = 2(w-3)$.
The proximal operator for $g(w) = |w|$ is the Soft-Thresholding operator:
$$\text{prox}_{\eta |\cdot|}(y) = S_\eta(y) = \text{sign}(y)\max(|y|-\eta, 0)$$

\textbf{Iteration 1:}
\begin{itemize}
    \item Gradient: $\nabla f(0) = 2(0-3) = -6$
    \item Forward step: $y_0 = w_0 - \eta \nabla f(w_0) = 0 - 0.1(-6) = 0.6$
    \item Proximal step: $w_1 = S_{0.1}(0.6) = \text{sign}(0.6)\max(0.6 - 0.1, 0) = 0.5$
    \item Objective: $F(0.5) = (0.5-3)^2 + |0.5| = 6.25 + 0.5 = 6.75$ (Down from $F(0) = 9$)
\end{itemize}

\textbf{Iteration 2:}
\begin{itemize}
    \item Gradient: $\nabla f(0.5) = 2(0.5-3) = -5$
    \item Forward step: $y_1 = w_1 - \eta \nabla f(w_1) = 0.5 - 0.1(-5) = 1.0$
    \item Proximal step: $w_2 = S_{0.1}(1.0) = \text{sign}(1.0)\max(1.0 - 0.1, 0) = 0.9$
    \item Objective: $F(0.9) = (0.9-3)^2 + |0.9| = 4.41 + 0.9 = 5.31$
\end{itemize}

\textbf{Iteration 3:}
\begin{itemize}
    \item Gradient: $\nabla f(0.9) = 2(0.9-3) = -4.2$
    \item Forward step: $y_2 = w_2 - \eta \nabla f(w_2) = 0.9 - 0.1(-4.2) = 1.32$
    \item Proximal step: $w_3 = S_{0.1}(1.32) = \text{sign}(1.32)\max(1.32 - 0.1, 0) = 1.22$
    \item Objective: $F(1.22) = (1.22-3)^2 + |1.22| = 3.1684 + 1.22 = 4.3884$
\end{itemize}

\section*{Assumptions}
To guarantee convergence, we make the following standard assumptions:
\begin{itemize}
    \item \textbf{A1 (Convexity and Smoothness of $f$):} The function $f: \mathbb{R}^d \to \mathbb{R}$ is convex and continuously differentiable. Furthermore, its gradient is $L$-Lipschitz continuous, meaning $\|\nabla f(x) - \nabla f(y)\| \le L \|x - y\|$ for all $x, y$.
    \item \textbf{A2 (Convexity of $g$):} The function $g: \mathbb{R}^d \to \mathbb{R} \cup \{+\infty\}$ is proper, closed, and convex.
    \item \textbf{A3 (Existence of Minimum):} The optimal set $\mathcal{X}^* = \arg\min_x F(x)$ is non-empty, and $F^* = F(x^*)$ is finite for $x^* \in \mathcal{X}^*$.
    \item \textbf{A4 (Step Size):} The step size satisfies $0 < \eta \le \frac{1}{L}$.
\end{itemize}

\section*{Main Convergence Theorem}
\textbf{Theorem:} Under assumptions A1-A4, the sequence $\{x_k\}$ generated by Proximal Gradient Descent satisfies the sublinear convergence rate:
$$F(x_K) - F(x^*) \le \frac{\|x_0 - x^*\|^2}{2\eta K}$$
where $x^*$ is any optimal solution.

\section*{Theoretical Properties}
\begin{itemize}
    \item \textbf{Stability via Non-expansiveness:} The proximal operator is firmly non-expansive, which ensures that applying the non-smooth penalty does not amplify distances between points. This guarantees strict Fejér monotonicity ($\|x_{k+1} - x^*\| \le \|x_k - x^*\|$).
    \item \textbf{Hyperparameter Sensitivity:} The algorithm is sensitive to the step size $\eta$. If $\eta > 2/L$, the algorithm may diverge. If $\eta$ is too small, convergence is unnecessarily slow. 
    \item \textbf{Comparison to Baselines:} Unlike subgradient descent which converges at a rate of $O(1/\sqrt{K})$ for general non-smooth convex functions, PGD achieves an $O(1/K)$ rate by exploiting the structure of $F(x)$, handling the non-smoothness analytically via the proximal operator.
\end{itemize}

\section*{Discussion}
Proximal Gradient Descent is a cornerstone algorithm in modern optimization, particularly for regularized machine learning problems like LASSO. It elegantly splits the objective, utilizing gradient information where the function is smooth, and analytical proximal mappings where it is not. The proof of its convergence heavily relies on the non-expansiveness of the proximal mapping and the co-coercivity of the gradient of the smooth component.

\section*{Preliminaries (Definitions and Lemmas)}

\textbf{Lemma 1 (Non-expansiveness of the Proximal Operator):} For any $x, y \in \mathbb{R}^d$, the proximal operator is non-expansive:
$$\|\text{prox}_{\eta g}(x) - \text{prox}_{\eta g}(y)\| \le \|x - y\|$$
\textit{Proof:} Let $u = \text{prox}_{\eta g}(x)$ and $v = \text{prox}_{\eta g}(y)$. By the optimality condition of the proximal operator, $x - u \in \eta \partial g(u)$ and $y - v \in \eta \partial g(v)$. By the monotonicity of the subdifferential for convex functions, $\langle (x - u) - (y - v), u - v \rangle \ge 0$. Rearranging yields $\langle x - y, u - v \rangle - \|u - v\|^2 \ge 0$. Applying the Cauchy-Schwarz inequality gives $\|u - v\|^2 \le \langle x - y, u - v \rangle \le \|x - y\| \|u - v\|$, which implies $\|u - v\| \le \|x - y\|$.

\textbf{Lemma 2 (Descent Lemma for Smooth Functions):} For a convex, $L$-smooth function $f$, for any $x, y$:
$$f(y) \le f(x) + \langle \nabla f(x), y - x \rangle + \frac{L}{2}\|y - x\|^2$$

\textbf{Lemma 3 (Gradient Co-coercivity):} For a convex, $L$-smooth function $f$ (Baillon-Haddad Theorem):
$$\langle \nabla f(x) - \nabla f(y), x - y \rangle \ge \frac{1}{L}\|\nabla f(x) - \nabla f(y)\|^2$$

\section*{Main Proof (Step-by-Step Derivation)}

We first utilize the non-expansiveness of the proximal operator to establish the stability of the iterations, and then derive the objective descent.

\noindent [Step 1] Define the forward-backward composite operator $T(x) = \text{prox}_{\eta g}(x - \eta \nabla f(x))$. The PGD algorithm is simply the fixed-point iteration $x_{k+1} = T(x_k)$.

\noindent [Step 2] We analyze the distance between gradient steps. For any $x, y$:
$$\|(x - \eta \nabla f(x)) - (y - \eta \nabla f(y))\|^2 = \|x - y\|^2 - 2\eta \langle \nabla f(x) - \nabla f(y), x - y \rangle + \eta^2 \|\nabla f(x) - \nabla f(y)\|^2$$

\noindent [Step 3] Applying Lemma 3 (Gradient Co-coercivity), we bound the cross-term:
$$-2\eta \langle \nabla f(x) - \nabla f(y), x - y \rangle \le -\frac{2\eta}{L} \|\nabla f(x) - \nabla f(y)\|^2$$

\noindent [Step 4] Substituting [Step 3] into [Step 2] yields:
$$\|(I - \eta \nabla f)(x) - (I - \eta \nabla f)(y)\|^2 \le \|x - y\|^2 + \eta\left(\eta - \frac{2}{L}\right)\|\nabla f(x) - \nabla f(y)\|^2$$

\noindent [Step 5] For $\eta \le \frac{1}{L} < \frac{2}{L}$, the term $\eta(\eta - 2/L) \le 0$. Thus, the gradient operator $I - \eta \nabla f$ is non-expansive:
$$\|(I - \eta \nabla f)(x) - (I - \eta \nabla f)(y)\| \le \|x - y\|$$

\noindent [Step 6] By Lemma 1, the proximal operator is non-expansive. Because the composition of two non-expansive operators is non-expansive, the forward-backward operator $T$ is non-expansive for $\eta \le 1/L$:
$$\|T(x) - T(y)\| \le \|x - y\|$$

\noindent [Step 7] Since $x^*$ minimizes $F(x)$, it satisfies $0 \in \nabla f(x^*) + \partial g(x^*)$, which implies $x^* = \text{prox}_{\eta g}(x^* - \eta \nabla f(x^*)) = T(x^*)$. Thus, $x^*$ is a fixed point of $T$.

\noindent [Step 8] Applying the non-expansiveness of $T$ from [Step 6] to $x_k$ and $x^*$, we obtain strict Fejér monotonicity:
$$\|x_{k+1} - x^*\| = \|T(x_k) - T(x^*)\| \le \|x_k - x^*\|$$
This guarantees the sequence of distances to the optimum is monotonically non-increasing.

\noindent [Step 9] Now we establish objective descent. By Lemma 2 (Smoothness of $f$):
$$f(x_{k+1}) \le f(x_k) + \langle \nabla f(x_k), x_{k+1} - x_k \rangle + \frac{L}{2} \|x_{k+1} - x_k\|^2$$

\noindent [Step 10] From the definition of the proximal operator, $x_{k+1} = \arg\min_z \{ g(z) + \frac{1}{2\eta}\|z - (x_k - \eta \nabla f(x_k))\|^2 \}$. The optimality condition requires that $0$ is in the subdifferential at $x_{k+1}$:
$$0 \in \partial g(x_{k+1}) + \frac{1}{\eta}(x_{k+1} - (x_k - \eta \nabla f(x_k)))$$

\noindent [Step 11] Rearranging the terms gives the subgradient of $g$ at $x_{k+1}$:
$$\frac{1}{\eta}(x_k - x_{k+1}) - \nabla f(x_k) \in \partial g(x_{k+1})$$

\noindent [Step 12] By the definition of convexity for $g$, for any arbitrary $x$:
$$g(x) \ge g(x_{k+1}) + \left\langle \frac{1}{\eta}(x_k - x_{k+1}) - \nabla f(x_k), x - x_{k+1} \right\rangle$$

\noindent [Step 13] Setting $x = x_k$ in [Step 12], we get:
$$g(x_k) \ge g(x_{k+1}) + \frac{1}{\eta}\|x_k - x_{k+1}\|^2 - \langle \nabla f(x_k), x_k - x_{k+1} \rangle$$
Which implies: $g(x_{k+1}) - g(x_k) \le -\frac{1}{\eta}\|x_k - x_{k+1}\|^2 - \langle \nabla f(x_k), x_{k+1} - x_k \rangle$.

\noindent [Step 14] Adding the inequality for $f$ from [Step 9] and $g$ from [Step 13]:
$$F(x_{k+1}) - F(x_k) \le \langle \nabla f(x_k), x_{k+1} - x_k \rangle + \frac{L}{2} \|x_{k+1} - x_k\|^2 - \frac{1}{\eta}\|x_k - x_{k+1}\|^2 - \langle \nabla f(x_k), x_{k+1} - x_k \rangle$$
Notice that the inner product terms containing the gradient perfectly cancel out.

\noindent [Step 15] This simplifies to a strictly decreasing objective (since $\eta \le 1/L \implies 1/\eta - L/2 \ge 1/(2\eta)$):
$$F(x_{k+1}) - F(x_k) \le -\left(\frac{1}{\eta} - \frac{L}{2}\right)\|x_{k+1} - x_k\|^2 \le -\frac{1}{2\eta}\|x_{k+1} - x_k\|^2 \le 0$$

\noindent [Step 16] Next, we bound the optimality gap. Set $x = x^*$ in the subgradient inequality from [Step 12]:
$$g(x^*) \ge g(x_{k+1}) + \left\langle \frac{1}{\eta}(x_k - x_{k+1}) - \nabla f(x_k), x^* - x_{k+1} \right\rangle$$

\noindent [Step 17] By the convexity of $f$, we have $f(x^*) \ge f(x_k) + \langle \nabla f(x_k), x^* - x_k \rangle$. Substituting $f(x_k)$ into [Step 9] yields:
$$f(x_{k+1}) \le f(x^*) + \langle \nabla f(x_k), x_k - x^* \rangle + \langle \nabla f(x_k), x_{k+1} - x_k \rangle + \frac{L}{2}\|x_{k+1} - x_k\|^2$$
$$f(x_{k+1}) - f(x^*) \le \langle \nabla f(x_k), x_{k+1} - x^* \rangle + \frac{L}{2}\|x_{k+1} - x_k\|^2$$

\noindent [Step 18] Adding the inequality for $g$ from [Step 16] and $f$ from [Step 17]:
$$F(x_{k+1}) - F(x^*) \le \left\langle \frac{1}{\eta}(x_k - x_{k+1}), x_{k+1} - x^* \right\rangle - \langle \nabla f(x_k), x_{k+1} - x^* \rangle + \langle \nabla f(x_k), x_{k+1} - x^* \rangle + \frac{L}{2}\|x_{k+1} - x_k\|^2$$

\noindent [Step 19] The gradient terms cancel out once again, leaving:
$$F(x_{k+1}) - F(x^*) \le \frac{1}{\eta}\langle x_k - x_{k+1}, x_{k+1} - x^* \rangle + \frac{L}{2}\|x_{k+1} - x_k\|^2$$

\noindent [Step 20] We apply the algebraic identity $2\langle a, b \rangle = \|a+b\|^2 - \|a\|^2 - \|b\|^2$ with $a = x_k - x_{k+1}$ and $b = x_{k+1} - x^*$. Note that $a+b = x_k - x^*$:
$$\langle x_k - x_{k+1}, x_{k+1} - x^* \rangle = \frac{1}{2}\left(\|x_k - x^*\|^2 - \|x_k - x_{k+1}\|^2 - \|x_{k+1} - x^*\|^2\right)$$

\noindent [Step 21] Substituting [Step 20] into [Step 19]:
$$F(x_{k+1}) - F(x^*) \le \frac{1}{2\eta}\left(\|x_k - x^*\|^2 - \|x_{k+1} - x^*\|^2\right) - \frac{1}{2\eta}\|x_k - x_{k+1}\|^2 + \frac{L}{2}\|x_{k+1} - x_k\|^2$$

\noindent [Step 22] Factoring the distance between consecutive points:
$$F(x_{k+1}) - F(x^*) \le \frac{1}{2\eta}\left(\|x_k - x^*\|^2 - \|x_{k+1} - x^*\|^2\right) - \left(\frac{1}{2\eta} - \frac{L}{2}\right)\|x_{k+1} - x_k\|^2$$

\section*{Rate Derivation (With Constants)}

\noindent [Step 23] Because we chose $\eta \le 1/L$, the term $\left(\frac{1}{2\eta} - \frac{L}{2}\right) \ge 0$. Therefore, we can drop the last term to find an upper bound:
$$F(x_{k+1}) - F(x^*) \le \frac{1}{2\eta}\left(\|x_k - x^*\|^2 - \|x_{k+1} - x^*\|^2\right)$$

\noindent [Step 24] We sum this inequality over $k = 0, 1, \dots, K-1$:
$$\sum_{k=0}^{K-1} (F(x_{k+1}) - F(x^*)) \le \frac{1}{2\eta} \sum_{k=0}^{K-1} \left(\|x_k - x^*\|^2 - \|x_{k+1} - x^*\|^2\right)$$

\noindent [Step 25] The right-hand side is a telescoping sum that collapses to the initial and final terms:
$$\sum_{k=0}^{K-1} (F(x_{k+1}) - F(x^*)) \le \frac{1}{2\eta} \left(\|x_0 - x^*\|^2 - \|x_K - x^*\|^2\right) \le \frac{1}{2\eta} \|x_0 - x^*\|^2$$

\noindent [Step 26] From [Step 15], we established that the sequence $F(x_k)$ is monotonically decreasing. Therefore, $F(x_K) \le F(x_{k+1})$ for all $k \le K-1$. This implies:
$$K(F(x_K) - F(x^*)) \le \sum_{k=0}^{K-1} (F(x_{k+1}) - F(x^*))$$

\noindent [Step 27] Combining [Step 25] and [Step 26], we obtain the final convergence rate:
$$F(x_K) - F(x^*) \le \frac{\|x_0 - x^*\|^2}{2\eta K}$$
This concludes the proof, demonstrating an $O(1/K)$ convergence rate.

\section*{Special Cases}
The Proximal Gradient Descent algorithm generalizes several fundamental algorithms depending on the choice of $f$ and $g$:
\begin{enumerate}
    \item \textbf{Standard Gradient Descent:} If $g(x) = 0$, the proximal operator becomes the identity mapping ($\text{prox}_{\eta g}(y) = y$), and the update reduces to $x_{k+1} = x_k - \eta \nabla f(x_k)$.
    \item \textbf{Projected Gradient Descent:} If $g(x) = I_C(x)$ (the indicator function of a convex set $C$, which is 0 if $x \in C$ and $+\infty$ otherwise), the proximal operator becomes the Euclidean projection operator onto $C$, $\Pi_C(y)$.
    \item \textbf{Iterative Shrinkage-Thresholding Algorithm (ISTA):} If $g(x) = \lambda \|x\|_1$, the proximal operator is the Soft-Thresholding operator, making this the standard algorithm for sparse recovery (e.g., LASSO).
    \item \textbf{Proximal Point Algorithm:} If $f(x) = 0$, there is no gradient step, and the algorithm relies entirely on the proximal operator, yielding $x_{k+1} = \text{prox}_{\eta g}(x_k)$.
\end{enumerate}

\end{document}